% Mục 3.1. Khối A (khoảng 0,5 trang). Người viết: Hiếu.
Khối A gồm năm bước chạy theo thứ tự cố định, mỗi bước là một hàm nhận ảnh và trả về ảnh, có thể tắt riêng (\path{code/src/enhance/}).
\begin{enumerate}[label=A\arabic*.]
  \item Ảnh được nắn phối cảnh \chuong{5, 2} dựa trên biên Canny (ngưỡng 50 và 150). Nhóm tìm tứ giác lớn nhất chiếm ít nhất 20\% diện tích ảnh rồi nắn nó bằng phép biến đổi phối cảnh với nội suy bậc ba, và giữ nguyên ảnh nếu không tìm được tờ giấy.
  \item Góc nghiêng được ước lượng \chuong{5} sau khi nhị phân hóa ảnh bằng Sauvola~\cite{sauvola2000binarization} (cửa sổ 31, $k=0{,}2$), bằng cách chọn góc $\theta^*$ sao cho phương sai của lược đồ chiếu ngang lớn nhất,
  \begin{equation}\label{eq:deskew}
    \theta^*=\arg\max_{\theta\in[-10^\circ,10^\circ]}\operatorname{Var}\bigl(P_\theta\bigr),\qquad P_\theta(y)=\sum_x b_\theta(x,y).
  \end{equation}
  Khi dòng chữ nằm ngang, mỗi hàng ảnh hoặc có nhiều mực hoặc gần như trống nên phương sai lớn nhất. Góc được tìm thô với bước $1^\circ$ rồi tinh lại với bước $0{,}25^\circ$, sau đó ảnh được xoay ngược lại theo góc này.
  \item Bóng được khử \chuong{7, 3} bằng cách ước lượng nền $\hat B$ qua phép đóng với phần tử cấu trúc vuông $81\times81$, rồi chia ảnh cho nền theo $g=\min\bigl(255,\;255\,f/\hat B\bigr)$ trên kênh V. Phần tử cấu trúc lớn hơn nhiều chiều cao chữ (trung vị 22~px) nên phép đóng làm mất nét chữ và chỉ còn lại ánh sáng nền.
  \item Nhiễu được giảm \chuong{3} bằng lọc trung vị $3\times3$, nhằm loại nhiễu hạt và nhiễu do nén JPEG.
  \item Tương phản được tăng \chuong{3} bằng CLAHE~\cite{zuiderveld1994clahe} trên kênh L, với giới hạn cắt 2,0 và lưới $8\times8$.
\end{enumerate}
Thí nghiệm độ bền dùng ba cấu hình là \texttt{none} (tắt hết), \texttt{light} (A2 và A3) và \texttt{full} (A1 đến A5).
\TODO{Hiếu: lý do chọn từng tham số, ghi khoảng đã thử nếu có và đưa bảng đầy đủ vào Phụ lục~\ref{app:params}}
